\documentclass[10pt,a4paper]{article}
\usepackage{amsmath,amssymb,amsthm,graphicx}
\usepackage{xfrac}
% Minimal editorial layout adapter for the pinned JEP source
% sources/jep-2022-195-contact-systolic/pinned-original-native.tex.
% The publisher document class (cedram) is not part of the retained material, and it
% is the class that used to define a handful of plain text abbreviations used by the
% author's prose. This adapter supplies exactly those abbreviations and nothing else:
% no mathematics, no notation, no theorem environment, no counter.
% HoTT/JEP retention convention: CC BY 4.0, attribution to the original authors.
\providecommand{\ie}{i.e.\ }
\providecommand{\eg}{e.g.\ }
\providecommand{\skpt}{\smallskip}
\usepackage[scr=rsfs,cal=euler]{mathalfa}
\let\mathds\mathbb
\let\epsilon\varepsilon
\newcommand{\Changel}{\mathcode`l="7160}
\newcommand{\Changelback}{\mathcode`l="716C}
\newcommand{\NoChangel}[1]{%
\expandafter\let\csname old\string#1\endcsname=#1
\let#1=\relax
\newcommand{#1}{\mathcode`l="716C\csname old\string#1\endcsname\mathcode`l="7160 }%
}

\NoChangel{\log}
\NoChangel{\ln}
\NoChangel{\lim}
\NoChangel{\limsup}
\NoChangel{\liminf}
\NoChangel{\varlimsup}
\NoChangel{\varliminf}
\NoChangel{\varinjlim}
\NoChangel{\varprojlim}

\newenvironment{enumeratei}
{\bgroup\def\theenumi{\roman{enumi}}\def\theenumii{\arabic{enumii}}\begin{enumerate}}
{\end{enumerate}\egroup}
\newcommand{\Psfrac}[2]{(\sfrac{#1}{#2})}
\newcommand{\psfrac}[2]{\sfrac{(#1)}{#2}}
\newcommand{\spfrac}[2]{\sfrac{#1}{(#2)}}
\newcommand\mto{\mathchoice{\longmapsto}{\mapsto}{\mapsto}{\mapsto}}
\newcommand\hto{\mathrel{\lhook\joinrel\to}}

\usepackage{graphicx}

\newtheorem{thm}{Theorem}[section]
\newtheorem{cor}[thm]{Corollary}
\newtheorem{lem}[thm]{Lemma}
\newtheorem{prop}[thm]{Proposition}
\newtheorem{quest}{Question}[section]
\newtheorem{conj}{Conjecture}[section]
\newtheorem{maintheorem}{Theorem}
\theoremstyle{definition}
\newtheorem{rem}[thm]{Remark}
\newtheorem{exm}[thm]{Example}
\newtheorem{ass}[thm]{Assumption}
\renewcommand*{\themaintheorem}{\Alph{maintheorem}}

\DeclareMathOperator{\area}{area}
\DeclareMathOperator{\lcm}{lcm}
\DeclareMathOperator{\const}{const}
\DeclareMathOperator{\Cal}{Cal}
\DeclareMathOperator{\Flux}{\mathrm{Flux}}
\newcommand{\id}{\mathrm{id}}
\newcommand{\Ham}{\mathrm{Ham}}
\newcommand{\std}{\mathrm{std}}
\newcommand{\vol}{\mathrm{vol}}
\newcommand{\fix}{\mathrm{fix}}
\newcommand{\pr}{\mathrm{pr}}
\newcommand{\Tan}{\mathrm{T}}
\newcommand{\N}{\mathds{N}}
\newcommand{\Z}{\mathds{Z}}
\newcommand{\R}{\mathds{R}}
\newcommand{\CC}{\mathds{C}}\let\C\CC
\newcommand{\UU}{\mathcal{U}}
\newcommand{\VV}{\mathcal{V}}

\let\eul e
\let\diff d
\title{Normalised almost-Besse contact forms: boundary behaviour and C1-closeness of the Reeb field}
\author{Alberto Abbondandolo, Christian Lange, Marco Mazzucchelli}
\date{doi:10.5802/jep.195; CC BY 4.0}
\begin{document}
\maketitle
\noindent\textit{Editorial note.} The counted statement is the author's labelled lemma `l:nearly\_Besse\_normalized` (source lines 1289-1310); the proof environment begins on the next line after \textbackslash{}end\{lem\} apart from one blank line, so the pairing is by adjacency.\par\medskip
\noindent\textit{Editorial note.} No context span is retained: the statement and proof are self-contained in the pinned author file, and the notation they use is either defined inside them or carried by the retained author macro block.\par\medskip
\noindent\textit{Editorial note.} The retained preamble is a minimal editorial adapter supplying only the plain text abbreviations that the publisher class used to define, followed by the native author macro block of the pinned source, excluding the publisher document class and its submission-date commands; the xfrac package is declared for the author's sfrac macros used by the retained text.\par\medskip
\noindent\textit{Editorial note.} No retained line refers to a label, citation or figure outside the retained span, so the package declares no external reference.\par\medskip
\noindent\textit{Editorial scope.} One original result of the pinned publisher article, counted once: the normalisation lemma for almost-Besse contact forms (author label l:nearly\_Besse\_normalized).\par\medskip
\noindent\textit{Allowed standard prerequisites (editorial).}\begin{itemize}
\item Almost-Besse contact form. A contact form is almost-Besse when the Reeb field is nowhere tangent to the contact structure and the linearised return map along each closed Reeb orbit has all eigenvalues real.. It is the hypothesis under which the four conclusions of the statement are derived.
\item Collar coordinates on the boundary. A collar neighbourhood A\_C x S\textasciicircum{}1 of a component C x S\textasciicircum{}1 of the boundary is identified with [0,rho) x S\textasciicircum{}1 x S\textasciicircum{}1, and the solid cylinder Y\_1 with B\_1 x S\textasciicircum{}1.. They are introduced at the start of the proof at source line 1313 to write the gluing map explicitly.
\item The map tilde-iota and the normalised form. The map tilde-iota(r,s,t) = (r e\textasciicircum{}\{2 pi i (q\_0 s - alpha\_1' t)\}, p\_0 s + alpha\_1 t) identifies the collar with the cylinder, and tilde-lambda\_0 is the displayed normalised contact form.. They are used at source lines 1313-1327 to compute the pull-back of nu and the vanishing of omega.
\end{itemize}\medskip
\begin{lem}\label{l:nearly_Besse_normalized}
\begin{enumeratei}

\item\label{l:nearly_Besse_normalizedi}
The pull-back of $\nu$ via the inclusion $\partial\Sigma\hto\Sigma$ satisfies
\[
\nu|_{\partial\Sigma}=\nu_0|_{\partial\Sigma}.
\]

\item\label{l:nearly_Besse_normalizedii}
The $2$-form $\omega$ vanishes at all points of $\partial\Sigma$, \ie
\[
\omega_z= 0,\qquad \forall z\in \partial\Sigma.
\]

\item\label{l:nearly_Besse_normalizediii}
The Reeb vector field $R_{\tilde\lambda}$, a priori only defined on the interior of $\Sigma\times S^1$, admits a smooth extension to $\Sigma\times S^1$ that is tangent to the boundary $\partial \Sigma\times S^1$.

\item\label{l:nearly_Besse_normalizediv}
For all $\epsilon>0$ there exists $\delta>0$ such that, if the above contact form $\lambda$ further satisfies $\|R_\lambda-R_{\lambda_0}\|_{C^2}<\delta$, then $\|R_{\tilde\lambda}-\partial_t\|_{C^1}<\epsilon$. Here, the norms are associated with arbitrary fixed Riemannian metrics.
\end{enumeratei}
\end{lem}

\begin{proof}
Let us consider a collar neighborhood $A_C\times S^1$ of a connected component $C\times S^1$ of $\partial\Sigma \times S^1$, and the solid cylinder neighborhood $Y_1\subset Y$ of $\gamma_1$. With the identifications $A_C\equiv[0,\rho)\times S^1$, $Y_1\equiv B_1\times S^1$, and $\gamma_1\equiv\{0\}\times S^1$, as in the previous subsection, the restriction of the map $\tilde\iota$ can be written as
\[
\tilde\iota|_{A_C\times S^1}:A_C\times S^1\to Y_1\equiv B_1\times S^1,
\qquad
\tilde\iota(r,s,t)=\big( r e^{2\pi(q_0s-\alpha_1't)i} , p_0s+\alpha_1t \big),
\]
and $\tilde\lambda_0|_{A_C\times S^1}$ can be written as
\begin{equation}
\label{e:pull_back_lambda_0}
\tilde\lambda_0|_{A_C\times S^1}
:=
\diff t
+
\Big(\frac{p_0}{\alpha_1} +\frac{2\pi}{\alpha_1} \big( q_0+\Psfrac{\alpha_1'}{\alpha_1}p_0 \big) r^2\Big)\diff s.
\end{equation}
The computations
\[
\tilde\iota_*\partial_t|_{C\times S^1} =R_{\lambda_0}=R_\lambda,
\qquad
\tilde\iota_*\partial_s|_{C\times S^1} = p_0\partial_s = \frac{p_0}{\alpha_1}\, R_{\lambda_0}= \frac{p_0}{\alpha_1}\, R_\lambda,
\]
imply
\[
\tilde\lambda(\partial_t)|_{C\times S^1}=\tilde\lambda_0(\partial_t)|_{C\times S^1} \equiv 1,
\qquad
\tilde\lambda(\partial_s)|_{C\times S^1} = \tilde\lambda_0(\partial_s)|_{C\times S^1} \equiv \frac{p_0}{\alpha_1},
\]
and
\[
\begin{aligned}
(\partial_t\,\lrcorner\,\diff\tilde\lambda)_z = \tilde\iota^*(R_{\lambda_0}\,\lrcorner\,\diff\lambda)_z &= \tilde\iota^*(R_\lambda\,\lrcorner\,\diff\lambda)_z \equiv 0,\\
(\partial_s\,\lrcorner\,\diff\tilde\lambda)_z
=
\frac{p_0}{\alpha_1}\,\tilde\iota^*(R_{\lambda_0}\,\lrcorner\,\diff\lambda)_z
&=
\frac{p_0}{\alpha_1}\,\tilde\iota^*(R_\lambda\,\lrcorner\,\diff\lambda)_z \equiv 0,
\end{aligned}
\qquad\forall z\in C\times S^1.
\]
These identities, together with $\nu_0=\tilde\lambda_0|_{\Sigma\times\{0\}}$ and $\nu=\tilde\lambda|_{\Sigma\times\{0\}}$, readily imply points~\eqref{l:nearly_Besse_normalizedi} and \eqref{l:nearly_Besse_normalizedii}

As for point \eqref{l:nearly_Besse_normalizediii}, let us write $R_\lambda|_{Y_1}$ in coordinates $(x+iy,s)=(re^{i\theta},s)$ as
\begin{align*}
R_\lambda|_{Y_1}
&= F_1\partial_x + F_2\partial_y + \partial_s\\
&= (F_1\cos(\theta)+F_2\sin(\theta))\,\partial_r + \frac1r\,(F_2\cos(\theta) - F_1\sin(\theta)) \partial_\theta + F_3\,\partial_s,
\end{align*}
for some smooth functions $F_j:Y_1\to\R$. Since $R_\lambda$ and $R_{\lambda_0}$ coincide along the closed Reeb orbit $\gamma_1$, we have
\[
F_1(0)=F_2(0)=0,\qquad F_3(0)=\alpha_1.
\]
Since
\[
\tilde\iota_*\partial_r=\partial_r,
\qquad
\tilde\iota_*\partial_s=2\pi q_0\,\partial_\theta + p_0\,\partial_s,
\qquad
\iota_*\partial_t= -2\pi\alpha_1'\,\partial_\theta + \alpha_1\,\partial_s,
\]
the Reeb vector fields $R_{\tilde\lambda}$ on the interior $\mathrm{int}(A_C)\times S^1$ is given by
\[
R_{\tilde\lambda}=f_1\,\partial_r + f_2\,\partial_s + f_3\,\partial_t,
\]
where, if we set
\[
\theta(s,t) :=q_0s-\alpha_1't,\qquad G_j(r,s,t):=\frac{F_j\circ\iota(r,s,t)}{r},\ \ j=1,2,
\]
the functions $f_j:\mathrm{int}(A_C)\times S^1\to\R$ are given by
\begin{align}
\label{e:f1}
f_1 & = \cos(\theta)F_1\circ\tilde\iota+\sin(\theta)F_2\circ\tilde\iota,\\
\label{e:f2}
f_2 & = \big( 2\pi\big(q_0+\Psfrac{\alpha_1'}{\alpha_1}p_0\big)\big)^{-1} \bigl( -\sin(\theta)G_1 +\cos(\theta)G_2 + 2\pi\Psfrac{\alpha_1'}{\alpha_1} F_3\circ\tilde\iota\bigr),\\
\label{e:f3}
f_3& = \frac1{\alpha_1}\big(F_3\circ\tilde\iota-p_0\,f_2\big).
\end{align}
Since $F_1\circ\tilde\iota|_{C\times S^1}=F_2\circ\tilde\iota|_{C\times S^1}\equiv0$, the functions $G_1$ and $G_2$ extend smoothly to the whole $A_C\times S^1$, and so do the functions $f_1,f_2,f_3$. Moreover, $f_1|_{C\times S^1}\equiv0$. This proves that $R_{\tilde\lambda}$ extends smoothly to a vector field on $\Sigma\times S^1$ that is tangent to the boundary $\partial\Sigma\times S^1$.

Finally, assume that $R_\lambda$ is $C^2$-close to $R_{\lambda_0}$. Away from any fixed neighborhood of $\partial\Sigma\times S^1$, $R_{\tilde\lambda}$ is $C^2$-close to $R_{\tilde\lambda_0}=\partial_t$. On $Y_1$, since $R_\lambda=-2\pi\alpha_1'\partial_\theta+\alpha_1\partial_s$, the functions
\[
F_1-2\pi\alpha_1' r\sin(\theta),\qquad
F_2+2\pi\alpha_1' r\cos(\theta),\qquad
F_3-\alpha_1
\]
are $C^2$-small and vanish at $\partial\Sigma\times S^1$. Therefore, the function
\[
-\sin(\theta)G_1+\cos(\theta)G_2+2\pi\Psfrac{\alpha_1'}{\alpha_1} F_3\circ\tilde\iota
\] is $C^1$-small. Equation \eqref{e:f1} implies that $f_1$ is $C^2$-small, and Equations \eqref{e:f2} and \eqref{e:f3} imply that $f_2$ and $f_3-1$ are $C^1$-small. Overall, we conclude that $R_{\tilde\lambda}$ is $C^1$-close to $R_{\tilde\lambda_0}=\partial_t$.
\end{proof}

\begin{thebibliography}{99}
\bibitem{Abbondandolo:2019tl} A. Abbondandolo \& G. Benedetti - “On the local systolic optimality of Zoll contact forms” , 2019
\bibitem{Abbondandolo:2018fb} A. Abbondandolo, B. Bramham, U. L. Hryniewicz \& P. A. S. Salomão - “Sharp systolic inequalities for Reeb flows on the three-sphere” , Invent. Math. 211 (2018) no. 2, p. 687-778
\bibitem{Abbondandolo:2019ta} A. Abbondandolo, B. Bramham, U. L. Hryniewicz \& P. A. S. Salomão - “Contact forms with large systolic ratio in dimension three” , Ann. Scuola Norm. Sup. Pisa Cl. Sci. (5) 19 (2019) no. 4, p. 1561-1582
\bibitem{Albach:2021wi} B. Albach \& H. Geiges - “Surfaces of section for Seifert fibrations” , Arnold Math. J. 7 (2021) no. 4, p. 573-597
\bibitem{Alvarez-Paiva:2014uq} J. C. Álvarez Paiva \& F. Balacheff - “Contact geometry and isosystolic inequalities” , Geom. Funct. Anal. 24 (2014) no. 2, p. 648-669
\bibitem{Banyaga:1978} A. Banyaga - “Sur la structure du groupe des difféomorphismes qui préservent une forme symplectique” , Comment. Math. Helv. 53 (1978) no. 2, p. 174-227
\bibitem{Banyaga:1997} A. Banyaga - The structure of classical diffeomorphism groups , Mathematics and its Appl. , vol. 400 , Kluwer Academic Publishers Group, Dordrecht, 1997
\bibitem{Besse:1978pr} A. L. Besse - Manifolds all of whose geodesics are closed , Ergeb. Math. Grenzgeb. , vol. 93 , Springer-Verlag, Berlin-New York, 1978
\bibitem{Benedetti:2021aa} G. Benedetti \& J. Kang - “A local contact systolic inequality in dimension three” , J. Eur. Math. Soc. (JEMS) 23 (2021) no. 3, p. 721-764
\bibitem{Bottkol:1980} M. Bottkol - “Bifurcation of periodic orbits on manifolds and Hamiltonian systems” , J. Differential Equations 37 (1980) no. 1, p. 12-22
\bibitem{Bialy:1994} M. Bialy \& L. Polterovich - “Geodesics of Hofer’s metric on the group of Hamiltonian diffeomorphisms” , Duke Math. J. 76 (1994) no. 1, p. 273-292
\bibitem{Calabi:1970} E. Calabi - “On the group of automorphisms of a symplectic manifold” , in Problems in analysis , Princeton University Press, Princeton, NJ, 1970, p. 1-26
\bibitem{Cristofaro-Gardiner:2020aa} D. Cristofaro-Gardiner \& M. Mazzucchelli - “The action spectrum characterizes closed contact 3-manifolds all of whose Reeb orbits are closed” , Comment. Math. Helv. 95 (2020) no. 3, p. 461-481
\bibitem{Croke:1994} C. B. Croke \& B. Kleiner - “Conjugacy and rigidity for manifolds with a parallel vector field” , J. Differential Geom. 39 (1994) no. 3, p. 659-680
\bibitem{Epstein:1972} D. B. A. Epstein - “Periodic flows on three-manifolds” , Ann. of Math. (2) 95 (1972), p. 66-82
\bibitem{Geiges:2020aa} H. Geiges - “What does a vector field know about volume?” , J. Fixed Point Theory Appl. 24 (2022) no. 2, article ID 23, 26 pages
\bibitem{Gromoll:1981kl} D. Gromoll \& K. Grove - “On metrics on $S^{2}$ all of whose geodesics are closed” , Invent. Math. 65 (1981/82) no. 1, p. 175-177
\bibitem{Ginzburg:1987lq} V. L. Ginzburg - “New generalizations of Poincaré’s geometric theorem” , Funktsional. Anal. i Prilozhen. 21 (1987) no. 2, p. 16-22
\bibitem{Giroux:2020} E. Giroux - “Ideal Liouville domains, a cool gadget” , J. Symplectic Geom. 18 (2020) no. 3, p. 769-790
\bibitem{Geiges:2018aa} H. Geiges \& C. Lange - “Seifert fibrations of lens spaces” , Abh. Math. Sem. Univ. Hamburg 88 (2018) no. 1, p. 1-22, Correction: Ibid. 91 (2021), no. 1, p. 145–150
\bibitem{Jankins:1983zm} M. Jankins \& W. D. Neumann - Lectures on Seifert manifolds , Brandeis Lect. Notes , vol. 2 , Brandeis University, Waltham, MA, 1983
\bibitem{Katok:1973mw} A. B. Katok - “Ergodic perturbations of degenerate integrable Hamiltonian systems” , Izv. Akad. Nauk SSSR Ser. Mat. 37 (1973), p. 539-576
\bibitem{Kegel:2021vo} M. Kegel \& C. Lange - “A Boothby-Wang theorem for Besse contact manifolds” , Arnold Math. J. 7 (2021) no. 2, p. 225-241
\bibitem{Lange:2020} C. Lange - “On metrics on 2-orbifolds all of whose geodesics are closed” , J. reine angew. Math. 758 (2020), p. 67-94
\bibitem{Lalonde:1995} F. Lalonde \& D. McDuff - “Hofer’s $L^\infty $-geometry: energy and stability of Hamiltonian flows. I, II” , Invent. Math. 122 (1995) no. 1, p. 1-33, 35–69, Erratum: Ibid. 123 (1996), no. 3, p. 613
\bibitem{Lisca:2004oz} P. Lisca \& G. Matić - “Transverse contact structures on Seifert 3-manifolds” , Algebraic Geom. Topol. 4 (2004), p. 1125-1144
\bibitem{Lange:2021} C. Lange \& T. Soethe - “Sharp systolic inequalities for rotationally symmetric $2$-orbifolds” , 2021
\bibitem{Mazzucchelli:2020aa} M. Mazzucchelli \& M. Radeschi - “On the structure of Besse convex contact spheres” , 2020
\bibitem{McDuff:1998uu} D. McDuff \& D. Salamon - Introduction to symplectic topology , Oxford Graduate Texts in Math. , Oxford University Press, Oxford, 2017
\bibitem{Orlik:1972ts} P. Orlik - Seifert manifolds , Lect. Notes in Math. , vol. 291 , Springer-Verlag, Berlin-New York, 1972
\bibitem{Radeschi:2017dz} M. Radeschi \& B. Wilking - “On the Berger conjecture for manifolds all of whose geodesics are closed” , Invent. Math. 210 (2017) no. 3, p. 911-962
\bibitem{Saglam:2018tb} M. Sağlam - “Contact forms with large systolic ratio in arbitrary dimensions” , Ann. Scuola Norm. Sup. Pisa Cl. Sci. (5) 22 (2021) no. 3, p. 1265-1308
\bibitem{Sullivan:1978} D. Sullivan - “A foliation of geodesics is characterized by having no ‘tangent homologies’” , J. Pure Appl. Algebra 13 (1978) no. 1, p. 101-104
\bibitem{Taubes:2007wi} C. H. Taubes - “The Seiberg-Witten equations and the Weinstein conjecture” , Geom. Topol. 11 (2007), p. 2117-2202
\bibitem{Wadsley:1975sp} A. W. Wadsley - “Geodesic foliations by circles” , J. Differential Geom. 10 (1975) no. 4, p. 541-549
\bibitem{Weinstein:1973} A. Weinstein - “Normal modes for nonlinear Hamiltonian systems” , Invent. Math. 20 (1973), p. 47-57
\bibitem{Ziller:1983rw} W. Ziller - “Geometry of the Katok examples” , Ergodic Theory Dynam. Systems 3 (1983) no. 1, p. 135-157
\end{thebibliography}
\end{document}
